% Victor: A Local-First Research Intelligence System with Epistemic Provenance
% Submission to AAAI-27 (41st AAAI Conference on Artificial Intelligence)
% Montréal, Canada, February 16--23, 2027
%
% COMPILATION:
%   1. Download the official AAAI-27 Author Kit from https://aaai.org/conference/aaai/aaai-27/
%   2. Place aaai27.sty (or aaai.sty) and aaai27.bst (or aaai.bst) in this directory
%   3. Run: pdflatex victor_aaai27 && bibtex victor_aaai27 && pdflatex victor_aaai27 && pdflatex victor_aaai27
%
% NOTE: Adjust the \usepackage{aaai} below to match the filename in the official kit
%       (e.g., \usepackage[submission]{aaai27} if the file is aaai27.sty).
%

% ==============================================================================
% PREAMBLE -- matches the official AAAI LaTeX template exactly
% ==============================================================================

\documentclass[letterpaper]{article}          % DO NOT CHANGE THIS
\usepackage[submission]{aaai}                 % DO NOT CHANGE THIS — use aaai27 if year-specific
\usepackage{times}                            % DO NOT CHANGE THIS
\usepackage{helvet}                           % DO NOT CHANGE THIS
\usepackage{courier}                          % DO NOT CHANGE THIS
\usepackage[hyphens]{url}                     % DO NOT CHANGE THIS
\usepackage{graphicx}                         % DO NOT CHANGE THIS
\urlstyle{rm}                                 % DO NOT CHANGE THIS
\def\UrlFont{\rm}                             % DO NOT CHANGE THIS
\usepackage{natbib}                           % DO NOT CHANGE THIS AND DO NOT ADD OPTIONS
\usepackage{caption}                          % DO NOT CHANGE THIS AND DO NOT ADD OPTIONS
\frenchspacing                                % DO NOT CHANGE THIS
\setlength{\pdfpagewidth}{8.5in}              % DO NOT CHANGE THIS
\setlength{\pdfpageheight}{11in}              % DO NOT CHANGE THIS

% -- Optional packages for algorithms/math (remove if not needed) ---------------
\usepackage{booktabs}                         % For table rules
\usepackage{amsmath}                          % For math environments
\usepackage{amssymb}                          % For extra math symbols
\usepackage{multirow}                         % For table spanning

% -- Theorem environments ------------------------------------------------------
\theoremstyle{plain}
\newtheorem{theorem}{Theorem}
\theoremstyle{definition}
\newtheorem{definition}[theorem]{Definition}
\theoremstyle{remark}
\newtheorem{remark}[theorem]{Remark}

% -- AAAI disallowed packages --------------------------------------------------
% DO NOT USE: authblk, balance, color (in text), CJK, float, flushend, fontenc,
% fullpage, geometry, grffile, hyperref, navigator, savetrees, setspace,
% stfloats, tabu, titlesec, tocbibind, ulem, wrapfig

% -- AAAI disallowed commands --------------------------------------------------
% DO NOT USE: \nocopyright, \addtolength, \balance, \baselinestretch, \clearpage,
% \columnsep, \newpage, \pagebreak, \pagestyle, \tiny

\setcounter{secnumdepth}{0}                   % Set to 1 or 2 if section numbers desired
\pdfinfo{/TemplateVersion (2027.1)}           % Update year as needed

% ==============================================================================
% TITLE & AUTHORS
% ==============================================================================

\title{Victor: A Local-First Research Intelligence System with Epistemic Provenance and NLI-Augmented Conflict Detection}

\author{Varish Sharma}
\affiliations{
    Independent Researcher \\
    varishsharma81106@gmail.com
}

\setcounter{theorem}{0}

\begin{document}

\maketitle

% ==============================================================================
% ABSTRACT
% ==============================================================================

\begin{abstract}
A researcher in the field does not carry a datacenter. They carry instruments, sensors, paper notes, and their own memory---fragmented tools with no unified understanding of what they have observed, what contradicts it, or what remains unknown. We present Victor, a research intelligence toolkit designed from the ground up for environments where cloud connectivity cannot be assumed: remote field sites, hazardous facilities, spacecraft, and submarines. Victor embeds a hierarchy of domain-specialist ML models---each trained for a specific sensor modality---that compress raw observations into structured inferences, which then feed a central reasoning model that builds compound understanding over time. Every piece of knowledge is stored in a masked memory graph: relational anchors between observations remain permanently live for routing and discovery, but the full content of any memory is hidden until the user explicitly calls it, ensuring sovereign control over sensitive research data. Victor operates at 90\% capability offline by default, carrying pre-trained models and dataset flows on local storage; internet adds capability but does not enable it. Security is two-tier: a defensive AI cybersecurity system guards each device instance independently, while a separate user-controlled switch determines what the system may publish to external networks. On the SciFact claim verification benchmark, the current software implementation achieves 71.0\% accuracy using Qwen2.5-7B and DeBERTa-v3-base NLI. Victor is open-source, CPU-only, and requires no cloud connectivity.
\end{abstract}

% ==============================================================================
% 1  INTRODUCTION
% ==============================================================================

\section{Introduction}

The distance between a physical observation and a researcher's reliable
understanding of it is enormous. Consider the motivating scenario: a person alone
in a hostile environment---a remote field site, a hazardous facility, or a
spacecraft---needs to understand their surroundings to survive. Their instruments
produce raw data. Their sensors capture observations. Their biological memory is
finite. And no single tool connects what was observed to what it means, what
contradicts it, or what remains unknown---without depending on cloud connectivity
that may not exist.

The real-world translation of this scenario is \emph{epistemic fragmentation}.
Today's researchers use instruments to observe, export tools to extract data,
notebooks to compute, reference managers to find papers, document viewers to read
them, note apps to record thoughts, and biological memory to hold it all together.
Every boundary between these tools destroys information. Every destroyed
connection must be reconstructed from memory months later, imperfectly, or not at
all. Experiments are run but their relationship to prior claims is never recorded.
Contradictions accumulate unseen. Gaps grow unnoticed.

We present Victor, a hierarchical multi-model research intelligence system that
addresses epistemic fragmentation by structuring knowledge as typed, provenance-
tracked artifacts from the moment of capture. Victor's architecture comprises
three layers: (1) domain-specialist ML models that compress raw sensor data into
structured inferences, (2) a central reasoning model that builds compound
understanding from accumulated inferences, and (3) a masked memory graph that
preserves relational structure while withholding content until the user explicitly
reveals it.

Victor is designed for deployment in environments where cloud connectivity cannot
be assumed. All inference, storage, and reasoning occur on-device. The system
carries pre-trained models locally and operates without internet access as the
default reliable state---90\% capable offline by design.

The key contributions of this paper are:

\begin{itemize}
    \item \textbf{Hierarchical multi-model reasoning architecture} --- a tiered
    system where domain-specialist models (vision, audio, extraction) feed a
    central reasoner through typed Epistemic Artifacts, with a strict trust
    hierarchy: MEASUREMENT $>$ EXPERIMENT $>$ LITERATURE $>$ NOTE $>$ HYPOTHESIS.
    \item \textbf{Masked memory graph with embedding-channel closure} --- a
    sovereignty primitive that suppresses node content from LLM answers while
    keeping graph structure live for routing, and deletes vector embeddings to
    close the membership-inference covert channel.
    \item \textbf{Two-tier sovereign security} --- a defensive AI cybersecurity
    system per device instance (tier 1) plus a separate user-controlled
    publication switch (tier 2), preventing any data from leaving the device
    without explicit authorization.
    \item \textbf{A bounded agent loop with automatic pre-answer review} --- a
    Retrieve $\rightarrow$ Review $\rightarrow$ Reason $\rightarrow$ Respond
    cycle with hard tool-call limits (5 calls), token ceilings (8,000 tokens),
    SHA-256 tool deduplication, and automatic conflict detection on retrieved
    claims before the first LLM call.
\end{itemize}

% ==============================================================================
% 2  RELATED WORK
% ==============================================================================

\section{Related Work}

\textbf{Retrieval-Augmented Generation.} \citet{lewis2020rag} established the
RAG paradigm: retrieve relevant documents, then generate an answer conditioned
on them. \citet{karpukhin2020dpr} introduced dense passage retrieval as the
standard retrieval backbone. Surveys \citep{gao2024ragsurvey} catalog the
rapidly expanding landscape. Victor extends RAG with typed epistemic
provenance---standard RAG has no notion of whether a retrieved chunk is a
hypothesis or a measurement. Victor's retrieval scoring blends cosine similarity
(0.7) with importance-weighted epistemic status (0.3), ensuring higher-tier
evidence ranks above lower-tier sources.

\textbf{Knowledge Graph QA.} KG-RAG systems \citep{pan2024kgrag} use knowledge
graphs for entity linking and relation traversal to improve retrieval.
\citet{ji2022graph} survey graph-based AI methods. Victor uses a knowledge graph
not for entity linking but for epistemic edge tracking: CONFIRMS, CONTRADICTS,
DERIVED\_FROM, and GROUNDED\_IN edges encode the evidential relationships between
claims, experiments, and observations. This typed edge layer enables automated
reasoning that flat knowledge graphs cannot support.

\textbf{NLI for Fact Verification.} \citet{thorne2018fever} and
\citet{wadden2020scifact} established NLI-based claim verification as a standard
evaluation task. \citet{devlin2019bert} and \citet{he2021deberta} provide the
encoder architectures underlying modern cross-encoders. The dominant approach
uses an LLM or fine-tuned classifier as the primary judge. Victor's contribution
is architectural: wiring the NLI cross-encoder as the primary direction judge
with the LLM as annotator-only, inverting the standard pipeline to eliminate
structured-output agreement bias where instruction-tuned LLMs default to
``confirm'' under uncertainty.

\textbf{Agentic RAG.} ReAct \citep{yao2023react}, Toolformer
\citep{schick2023toolformer}, and WebGPT \citep{nakano2021webgpt} introduce
tool-use loops where an LLM iteratively calls retrieval and reasoning tools.
\citet{wei2022cot} demonstrates chain-of-thought reasoning for complex tasks.
Victor bounds this loop explicitly with a hard tool-call limit (5 calls), a soft
token ceiling (8,000 tokens with forced-answer fallback), and an automatic Review
step that runs the conflict detector on retrieved claims with mixed evidence
before the first LLM call. SHA-256 tool deduplication prevents redundant calls.

\textbf{Memory-Augmented Systems.} \citet{graves2016hybrid} and
\citet{sukhbaatar2015memory} established differentiable external memory for
neural networks. Recent systems like MemGPT, Zep, and Mem0 bring vector-based
memory to LLMs but support deletion, not selective content masking while keeping
graph structure alive. Victor's approach is inspired by scope-before-routing
\citep{shardmemo2025} and retrieval-phase enforcement \citep{sdrag2025}---but
adapted for user-controlled revelation on a local device, not privacy from remote
LLMs.

\textbf{Personal Knowledge Management.} Tools like Obsidian, Roam Research, and
Logseq capture \emph{what} the user wrote but not \emph{how confident} they
should be in it. Victor differs by enforcing epistemic typing (every node carries
an epistemic status), automated conflict detection (experiments can confirm or
contradict claims), and evidence weighting (confidence scores aggregate from
evidence edges).

% ==============================================================================
% 3  SYSTEM DESIGN
% ==============================================================================

\section{System Design}

\subsection{Hierarchical Architecture}

Victor's architecture follows a three-layer hierarchy:

\begin{enumerate}
    \item \textbf{Domain-specialist layer:} Per-modality ML models compress raw
    sensor data into structured inferences. The vision classifier uses
    CLIP-ViT-B/32 \citep{radford2021clip} for zero-shot image classification.
    The audio classifier uses librosa feature extraction (MFCC, spectral
    centroid, zero-crossing rate, RMS energy) with heuristic scoring. File
    extractors parse PDF, Markdown, and JSON experiment logs.
    \item \textbf{Central reasoner layer:} A single reasoning model
    (Qwen2.5-7B, approximately 5 GB in Q4\_K\_M quantization) receives
    structured inferences from all domain specialists and builds compound
    understanding through a bounded agent loop.
    \item \textbf{Masked memory graph layer:} A NetworkX-directed graph with
    8 node types and 12 typed edge types stores all knowledge with full
    provenance. Content masking is enforced at the retrieval layer.
\end{enumerate}

\subsection{The Epistemic Artifact}

Victor's fundamental unit of knowledge is the Epistemic Artifact---a graph node
with six mandatory properties:

\begin{enumerate}
    \item \textbf{Origin} --- provenance chain: sensor type and ID, source
    document, or extraction model.
    \item \textbf{Type} --- one of five epistemic tiers forming a trust hierarchy:
    \begin{itemize}
        \item MEASUREMENT (weight 1.00): directly observed data point
        \item EXPERIMENT (weight 0.85): controlled test result
        \item LITERATURE (weight 0.70): claim from a published paper
        \item NOTE (weight 0.50): personal annotation or summary
        \item HYPOTHESIS (weight 0.25): speculative, unverified
    \end{itemize}
    \item \textbf{Confidence} --- 0.0--1.0, computed from evidence edge aggregation.
    \item \textbf{Contradictions} --- set of experiment nodes with CONTRADICTS
    edges to this claim.
    \item \textbf{Gaps} --- what evidence is missing before this claim can be
    trusted.
    \item \textbf{Age} --- timestamp plus temporal decay rate for importance
    weighting.
\end{enumerate}

The hierarchy is enforced at graph construction time: the node type determines
the default epistemic status, which determines the weight used in all downstream
retrieval scoring and conflict resolution. A MEASUREMENT node from a spectrometer
and a HYPOTHESIS node from casual observation are structurally distinct---the
system cannot confuse them.

\subsection{Knowledge Graph}

Victor uses a directed graph with 8 node types (PAPER, CLAIM, EXPERIMENT, NOTE,
CONCEPT, CONVERSATION, COMPUTED, OBSERVATION) and 12 typed edge types (MAKES\_CLAIM,
TESTED\_BY, CONFIRMS, CONTRADICTS, INCONCLUSIVE, REFERENCES, IMPLEMENTS,
RELATED\_TO, CONTAINS, GROUNDED\_IN, DERIVED\_FROM, SIMILAR\_TO). The graph is
serialized as JSON via NetworkX with atomic save semantics (write to \texttt{.tmp},
then \texttt{os.replace()}) to survive crashes. Every node carries semantic
embeddings computed at creation time for cosine-similarity retrieval.

\subsection{Ingestion Pipeline}

The ingestion layer uses a plugin architecture. Adding a new source type requires
implementing \texttt{extensions} and \texttt{ingest()}---the core controller
dispatches by file extension through a registry without modification. Three built-in
ingestors:

\begin{itemize}
    \item \textbf{PDFIngestor} --- extracts text via PyMuPDF/pdfplumber, chunks
    via \texttt{TextChunker}, optionally extracts claims via LLM-based
    \texttt{ClaimExtractor}.
    \item \textbf{MarkdownIngestor} --- same pipeline for \texttt{.md}/\texttt{.txt}/
    \texttt{.rst} files, accepts a \texttt{source\_type} override.
    \item \textbf{ExperimentIngestor} --- ingests structured JSON experiment logs,
    auto-links to existing claims via cosine similarity threshold.
\end{itemize}

Claims are extracted per-chunk (not on a hard-truncated document prefix) to avoid
losing information from later sections. Extracted claims are deduplicated
case-insensitively across chunks and stored with source provenance.

\subsection{Sensor Observation Pipeline}

Victor supports live observation from camera and microphone sources, converting
raw sensor data directly into MEASUREMENT-tier graph nodes.

\begin{itemize}
    \item \textbf{CameraIngestor} --- captures one frame from a webcam or loads
    an existing image file. Classification uses a CLIP-ViT-B/32 zero-shot model
    \citep{radford2021clip} against user-provided candidate labels. The image is
    saved for provenance and an \texttt{ObservationResult} with all six Epistemic
    Artifact properties is created.
    \item \textbf{AudioIngestor} --- records from microphone for a configurable
    duration or loads an existing audio file. Classification uses librosa-based
    feature extraction (MFCCs, spectral centroid, zero-crossing rate, RMS energy)
    with heuristic scoring. The audio clip is saved for provenance.
\end{itemize}

Both ingestors accept \texttt{extra\_gaps} parameters for researcher-specified
verification requirements. Mock classifiers are available for testing without
GPU or model download.

\subsection{Retrieval Architecture}

\textbf{HybridRetriever} performs two-stage retrieval with blended scoring:

\[
\text{score} = 0.7 \times \text{cosine\_similarity} + 0.3 \times \text{importance\_score}
\]

Stage 1 queries the ChromaDB vector store for top-k similar chunks. Stage 2
performs semantic graph query on node embeddings plus 2-hop neighbor expansion
via typed edges. Masked nodes participate in graph routing but their content is
withheld from retrieval results---emitting anchor chunks that preserve referential
presence without content exposure.

\textbf{AgenticRetriever} extends HybridRetriever with a bounded tool-use loop:

\begin{itemize}
    \item \textbf{Retrieve} --- hybrid vector + graph retrieval
    \item \textbf{Review} --- automatic conflict detection on retrieved claims
    with mixed evidence, before the first LLM call
    \item \textbf{Reason} --- LLM calls reasoning tools (judge\_conflict,
    summarize\_claim\_evidence, check\_claims\_in\_context, find\_research\_gaps,
    search\_nodes\_by\_text, find\_related\_nodes, get\_node\_details,
    get\_claims\_for\_source)
    \item \textbf{Respond} --- LLM generates the final answer conditioned on all
    gathered evidence
\end{itemize}

Hard limits: maximum 5 tool calls, 8,000 token ceiling with forced-answer
fallback. Tool deduplication by SHA-256 hash prevents redundant calls.

\subsection{Conflict Detection}

The \textbf{ConflictDetector} determines whether an experiment confirms,
contradicts, or is inconclusive with respect to a claim. It operates in two modes:

\textbf{Standard mode (LLM + NLI cross-check):} The LLM (Qwen2.5-7B) makes the
direction decision via structured JSON output. An NLI cross-encoder
(\texttt{cross-encoder/nli-deberta-v3-base}) runs afterward. If the LLM predicts
confirm/inconclusive but the NLI predicts contradiction above a confidence
threshold (default 0.5), the judgment is overridden.

\textbf{NLI-primary mode:} The NLI cross-encoder decides the direction. The LLM
is called only for a free-form reasoning trace with no JSON schema, eliminating
structured-output agreement bias.

The \textbf{NLICrossEncoder} wrapper handles label-order ambiguity across model
families. Victor auto-detects the output order from \texttt{model.config.id2label},
trying three attribute paths for robustness. A pre-benchmark smoke test validates
the mapping before any evaluation run.

\subsection{Reasoning Layer}

\textbf{EvidenceWeighter} walks inbound edges to a CLAIM node and computes
confidence as $n_{\text{confirms}} / (n_{\text{confirms}} + n_{\text{contradicts}} +
n_{\text{inconclusive}})$. \textbf{GapDetector} surfaces CLAIM nodes with no
confirming experiments, sorted by importance. \textbf{PatternDetector} performs
cross-gap analysis and auto-injects insights into answers when query embedding
similarity to a priority gap exceeds 0.65.

\subsection{Sovereignty and Security}

Victor implements two orthogonal per-node control flags:

\begin{itemize}
    \item \textbf{Freeze} prevents overwrite of node content while keeping it
    readable.
    \item \textbf{Mask} suppresses node content from LLM answers while keeping
    graph structure live for routing, gap detection, and conflict detection.
    Masked nodes emit anchor chunks preserving referential presence.
\end{itemize}

Masking also deletes the node's embedding chunks from ChromaDB, closing the
membership-inference covert channel where an adversary could triangulate masked
content via cosine similarity queries. Every operation is logged in a tamper-
evident JSONL audit trail with caller attribution.

Masking is enforced not only in the retrieval layer but in every tool in the
agent loop---all eight tools check the \texttt{is\_masked} flag before returning
node content. This closes a critical bypass vector where an LLM receiving a
masked anchor chunk (``Use victor unmask {node\_id} to reveal'') could call
\texttt{get\_node\_details(node\_id)} to read the full content without user
authorization.

\subsection{Computation Kernels}

The \textbf{StatisticsKernel} computes descriptive statistics (mean, variance,
standard deviation, confidence intervals) and Cohen's d effect size over node
values. It requires an explicit \texttt{metric\_name} parameter---never averaging
all numeric fields, preventing conflation of distinct metrics (e.g., precision +
recall + f1). Outputs are stored as COMPUTED nodes with DERIVED\_FROM edges and
a SHA-256 reproducibility hash.

% ==============================================================================
% 4  IMPLEMENTATION
% ==============================================================================

\section{Implementation}

Victor is implemented in Python (v5.3.0) as a package of 16 business logic
modules totaling approximately 6,000 lines of code (18 total \texttt{.py} files
including \texttt{\_\_init\_\_.py} imports and a standalone benchmark script).
The system is organized around a central \texttt{Victor} controller class that
orchestrates the following components:

\begin{itemize}
    \item \textbf{LocalEmbedder} --- wraps \texttt{sentence-transformers} with a
    fallback chain. Primary model is nomic-embed-text-v1.5; falls back to
    all-MiniLM-L6-v2.
    \item \textbf{VictorVectorStore} --- ChromaDB-backed vector store for
    chunk-level semantic retrieval with node-level deletion support for masked
    memory.
    \item \textbf{VictorKnowledgeGraph} --- NetworkX-based directed graph
    serialized to JSON with atomic save semantics. Stores typed nodes and edges
    with all epistemic metadata.
    \item \textbf{InferenceClient} --- provider-agnostic abstraction over Ollama
    (Qwen2.5-7B, primary), OpenAI API, and a mock backend for testing.
    \item \textbf{ToolRegistry} --- 8 reasoning tools registered for the agent
    loop, with SHA-256 deduplication and masking enforcement in every tool.
    \item \textbf{StatisticsKernel} --- NumPy/SciPy-based computation with
    explicit metric selection.
    \item \textbf{Sensor pipeline} --- CameraIngestor (CLIP-ViT-B/32 classifier),
    AudioIngestor (librosa feature classifier), with mock classifiers for testing.
\end{itemize}

Dependencies include \texttt{chromadb}, \texttt{sentence-transformers},
\texttt{networkx}, \texttt{ollama}, \texttt{click}, and \texttt{rich}. The system
runs entirely on CPU---the largest evaluated configuration (Qwen2.5-7B Q4\_K\_M,
approximately 5 GB) operates on consumer hardware with 32 GB RAM.

The CLI provides the complete interface via \texttt{click} and \texttt{rich}:
\texttt{victor ingest}, \texttt{victor ask}, \texttt{victor gaps},
\texttt{victor observe}, \texttt{victor inspect}, \texttt{victor stats},
\texttt{victor mask}, \texttt{victor unmask}, \texttt{victor freeze},
\texttt{victor unfreeze}, \texttt{victor export}, \texttt{victor delete},
and computation kernel commands. No web API or graph visualization is
implemented---these are out of scope for the current version.

The complete source code is publicly available under an open-source license. The test suite contains 196 tests across 12 test files covering the agent loop,
audio pipeline, CLI imports, masked memory, masking enforcement, metric
conflation, NLI cross-encoder, SciFact benchmark, sensor pipeline, statistics
kernel, tool mask enforcement, and audit logging.

% ==============================================================================
% 5  EVALUATION
% ==============================================================================

\section{Evaluation}

\subsection{Task and Dataset}

We evaluate Victor's ConflictDetector on the SciFact benchmark
\citep{wadden2020scifact}: given a (claim, abstract) pair, classify the
relationship as confirm, contradict, or inconclusive. The dev set contains 300
claims with evidence-linked abstracts. All experiments use n=100 samples unless
otherwise noted.

\subsection{Results}

\begin{table}[t]
\centering
\caption{Conflict detection accuracy on SciFact across five configurations.
Run 5 (DeBERTa cross-check) achieves the best result at 71.0\%.}
\label{tab:results}
\begin{tabular}{@{}lccccc@{}}
\toprule
Run & Mode & n & Accuracy & Confirm & Contradict \\
     &      &   &           & Recall   & Recall \\
\midrule
1 & Pure LLM, free-form & 50 & 68.0\% & 86.0\% & 27.0\% \\
2 & Structured JSON     & 50 & 58.0\% & 63.0\% & 47.0\% \\
3 & Qwen2.5-7B, no NLI & 100& 59.0\% & 61.4\% & 53.3\% \\
4 & DeBERTa, NLI primary& 100& 8.0\%  & 7.1\%  & 10.0\% \\
5 & DeBERTa, NLI cross-check & 100 & \textbf{71.0\%} & 68.6\% & 76.7\% \\
\bottomrule
\end{tabular}
\end{table}

Table~\ref{tab:results} summarizes the results across five configurations.

\textbf{Run 1 $\rightarrow$ Run 2 regression.} Forcing Qwen2.5-7B to emit valid
JSON against an enum-constrained schema suppresses reasoning quality. The model
routes to the lowest-perplexity token satisfying the schema---``confirm'' is
shorter, more common in RLHF data, and schema-valid. This produces a 10-point
accuracy drop.

\textbf{Run 4 failure (label-order bug).} The NLI label order constant had
indices 1 and 2 swapped. Every entailment prediction mapped to ``inconclusive,''
then overridden to ``contradict'' by primary-mode logic. The fix auto-detects the
order from \texttt{model.config.id2label} and adds a pre-benchmark smoke test.

\textbf{Run 5 (current best).} Cross-check mode is robust to the label-order bug
because it only uses index 0 (contradiction) for the override decision. The NLI
cross-encoder catches agreement bias on approximately 8 of 100 samples, adding
approximately 12 points over the pure-LLM baseline.

\subsection{Analysis}

The 71.0\% result represents the combined ceiling of Qwen2.5-7B (pure
approximately 59\%) plus DeBERTa-v3-base NLI cross-check. The NLI override count
(5--8 of 100) indicates the cross-encoder fires infrequently---the LLM performs
92--95\% of the work. Threshold tuning across values 0.3--0.6 did not improve on
71.0\%, confirming that override frequency, not threshold, is the binding
constraint.

Paths to 85\% accuracy include: (a) a scientifically fine-tuned NLI model
(SciFact-trained DeBERTa), (b) a larger LLM, (c) self-consistency voting, or
(d) LoRA fine-tuning on the SciFact training split.

% ==============================================================================
% 6  DISCUSSION
% ==============================================================================

\section{Discussion}

\textbf{Agreement bias in LLM-based NLI.} Instruction-tuned LLMs performing
zero-shot NLI systematically default to ``entailment'' (confirm) when uncertain.
This is training-distribution driven---RLHF data contains far more agreement than
contradiction. Forcing structured JSON output amplifies this by routing to the
schema-satisfying token with lowest perplexity. Victor's NLI-primary architecture
addresses this by removing the direction decision from the LLM entirely, though
the cross-encoder alone lacks the reasoning depth for complex scientific abstracts.
The cross-check architecture---LLM decides, NLI catches agreement bias---currently
achieves the best tradeoff.

\textbf{Label-order ambiguity.} NLI model families output label probabilities in
different orders. Victor resolves this through auto-detection from the model
config, an explicit \texttt{label\_order} constructor parameter, and a
pre-benchmark smoke test. This infrastructure caught a real failure (ModernCE's
inverted order) that would have produced 35\% accuracy without intervention.

\textbf{Embedding covert channel.} Setting a masking flag on a graph node
suppresses content from LLM answers but leaves vector embeddings queryable in
ChromaDB. Victor closes this channel by deleting the node's chunks from ChromaDB
at mask time. This is necessary because the vector store operates independently
of the graph's masking state.

\textbf{Masking enforcement across the entire tool surface.} A masked node in
the retrieval layer is insufficient---the LLM can call \texttt{get\_node\_details}
or \texttt{search\_nodes\_by\_text} to read content without passing through the
retrieval layer. Victor enforces masking in all eight tools of the agent loop,
ensuring consistent content suppression regardless of the access path.

% ==============================================================================
% 7  LIMITATIONS
% ==============================================================================

\section{Limitations}

We acknowledge the following limitations:

\begin{enumerate}
    \item \textbf{CPU-only inference limits model size.} Qwen2.5-7B is the
    largest model that runs comfortably on consumer hardware. Larger models or
    fine-tuned variants would likely improve accuracy but require GPU
    infrastructure.
    \item \textbf{Audio pipeline uses heuristic classification.} The
    AudioIngestor employs MFCC-based features, not deep neural network
    classifiers as originally planned. Classification accuracy for real-world
    audio events is unvalidated.
    \item \textbf{No web API or graph visualization.} The system is CLI-only.
    Interactive graph exploration and web-based interfaces are out of scope for
    the current version.
    \item \textbf{SciFact target not yet met.} The 85\% accuracy threshold
    remains open. The current best of 71.0\% represents the combined ceiling of
    the current model configuration.
    \item \textbf{Evaluation on benchmark pairs only.} No real research data has
    been ingested and evaluated end-to-end. Validation has been on claim-abstract
    pairs, not full document ingestion and gap detection workflows.
    \item \textbf{Single-seed evaluation at n=100.} The 71.0\% result has wide
    confidence intervals. A full 300-sample, two-seed evaluation is required
    before the result can be considered stable.
    \item \textbf{Physical device not yet fabricated.} The current implementation
    is a software prototype. The physical form factor (integrated sensor suite,
    onboard compute, sovereign network interface) is the target for future work.
\end{enumerate}

% ==============================================================================
% 8  CONCLUSION
% ==============================================================================

\section{Conclusion}

Victor demonstrates that hierarchical multi-model reasoning---domain specialists
feeding a central reasoner through typed Epistemic Artifacts---can be implemented
as a coherent local-first system without cloud dependency. The masked memory
graph with embedding-channel closure establishes that content suppression and
graph routing can coexist, and that the vector embedding covert channel must be
explicitly closed, not just flagged.

The NLI-primary architecture identifies and addresses structured-output agreement
bias in LLM-based fact verification. While the current cross-encoder lacks
reasoning depth for primary-mode deployment, the cross-check configuration
achieves 71.0\% on SciFact with clear paths to improvement through fine-tuned
NLI models, self-consistency voting, or larger LLMs.

The two-tier sovereign security model---defensive AI per device instance plus
user-controlled publication---provides a principled foundation for research
intelligence systems operating in sensitive or disconnected environments.

Future work targets the 85\% SciFact threshold via fine-tuned NLI models,
followed by MCP server integration for distribution through Claude Desktop and
Cursor, and ultimately the physical device platform with integrated multi-sensor
fusion.

% ==============================================================================
% REFERENCES
% ==============================================================================

\bibliography{references}

\end{document}
